% main_report81.tex
% SAMBP Technical Report TR-81/2028
% GMD/GIC Protection
\documentclass[11pt, a4paper]{article}

\usepackage[T1]{fontenc}
\usepackage[utf8]{inputenc}
\usepackage[margin=25mm]{geometry}
\usepackage{amsmath, amssymb, amsthm}
\usepackage{booktabs, array, multirow, longtable}
\usepackage{graphicx}
\usepackage[table]{xcolor}
\usepackage[hidelinks]{hyperref}
\usepackage{pgfplots}
\pgfplotsset{compat=1.18}
\usepackage{tikz}
\usetikzlibrary{arrows.meta, positioning, calc, fit, backgrounds,
                shapes.geometric, shapes.misc}
\usepackage{caption}
\usepackage{subcaption}
\usepackage{parskip}
\usepackage{siunitx}
\usepackage{pifont}
\usepackage{cite}

\definecolor{sambpblue}{RGB}{0,70,140}
\definecolor{okgreen}{RGB}{0,130,60}
\definecolor{alertred}{RGB}{180,0,0}
\definecolor{warnoran}{RGB}{200,100,0}

\newcommand{\pu}{\,\text{pu}}
\newcommand{\ms}{\,\text{ms}}
\newcommand{\cmark}{\ding{51}}
\newcommand{\xmark}{\ding{55}}

\newtheorem{finding}{Finding}
\newtheorem{remark}{Remark}

\title{\textbf{\textsf{\color{sambpblue}SAMBP Technical Report TR-81/2028}}\\[6pt]
  \Large Geomagnetic Disturbance and GIC Protection\\[4pt]
  \normalsize Lehtinen-Pirjola GIC Flow Model, IEC~60076-7 Transformer
  Thermal Model, 5-Bus 400\,kV Network: 11 Tests PASS}
\author{SAMBP Research Group\\[2pt]
  \small Smart Adaptive Multi-Bus Protection Research, IIT Madras}
\date{April 2028\quad\textbar\quad Chapter~3 (Protection Challenges)\quad\textbar\quad
  Prerequisite: TR-67 (HIL)}

\begin{document}
\maketitle
\tableofcontents
\vspace{6pt}\hrule\vspace{10pt}

\section{Background and Objectives}
\label{sec:background}

\subsection{Motivation}

Geomagnetic Disturbances (GMD) caused by solar storms drive quasi-DC
Geomagnetically Induced Currents (GIC) through transmission networks.
GIC saturates power transformers, causing half-cycle saturation, increased
reactive power demand, harmonic injection, and potentially catastrophic
thermal damage~\cite{Gaunt2007}. The Quebec 1989 blackout (6 hours, $9\,\text{GW}$
loss) and the Hydro-Quebec transformer damage are canonical examples.

The SAMBP framework adds GMD/GIC protection as a special protection scheme
(SPS) separate from the conventional fault-protection functions.

\subsection{Objectives}

\begin{enumerate}
  \item Implement the Lehtinen-Pirjola (LP) GIC flow model.
  \item Implement the IEC~60076-7:2018 transformer thermal model with GIC.
  \item Define alarm levels: NORMAL, ELEVATED, CRITICAL, TRIP.
  \item Validate on a 5-bus 400\,kV network with 11 unit tests.
\end{enumerate}

\section{Theory and Methodology}
\label{sec:theory}

\subsection{Lehtinen-Pirjola GIC Flow Model}

The LP model~\cite{Lehtinen1985} treats the transmission network as a
resistive circuit driven by the geoelectric field $\mathbf{E}$
(induced by the time-varying geomagnetic field). The network equations are:

\begin{equation}
  (\mathbf{Y}_n + \mathbf{Y}_e)\,\mathbf{V} = \mathbf{J}_e
  \label{eq:lp}
\end{equation}

\noindent where $\mathbf{Y}_n$ is the network admittance matrix
(DC resistance of lines and transformer windings),
$\mathbf{Y}_e$ is the earth admittance matrix (earthing resistances),
$\mathbf{V}$ is the vector of node voltages,
and $\mathbf{J}_e = \mathbf{Y}_e^{(ij)} \cdot \int_{ij} \mathbf{E}\,dl$
is the geoelectric source current vector (line integral of $\mathbf{E}$
along each transmission line).

The neutral GIC at bus $k$ is:
$I_{\mathrm{GIC},k} = Y_{e,kk} V_k$ (quasi-DC neutral current through transformer).

\paragraph{3-bus LP benchmark.}
With $R_{\mathrm{line}} = R_{\mathrm{earth}} = 2\,\Omega$ and
$E_x = 1\,\text{V/km}$:
Bus 1: $I_{\mathrm{GIC}} = 25.0\,\text{A}$;
Bus 2: $25.0\,\text{A}$; Bus 3: $-50.0\,\text{A}$.
KCL check: $\sum I_{\mathrm{GIC}} = 0.000\,\text{A}$. \cmark

\paragraph{5-bus 400\,kV scaling.}
GIC scales linearly with E-field intensity:
$I_{\mathrm{GIC,max}} = 57.7\,\text{A}$ at $E = 1\,\text{V/km}$;
$= 577\,\text{A}$ at $E = 10\,\text{V/km}$ (extreme solar storm).

\subsection{Transformer Thermal Model}

The IEC~60076-7:2018 thermal model (Thermal Model~1)~\cite{IEC60076}:

\begin{align}
  \tau_{\mathrm{TO}}\frac{d\Delta\Theta_{\mathrm{TO}}}{dt}
    &= \Delta\Theta_{\mathrm{TO,R}}
    \left[\frac{K_{\mathrm{eff}}^2 + R}{1 + R}\right]^{n_{\mathrm{TO}}}
    - \Delta\Theta_{\mathrm{TO}}
  \label{eq:topol}\\
  \tau_{\mathrm{HS}}\frac{d\Delta\Theta_{\mathrm{HS}}}{dt}
    &= \Delta\Theta_{\mathrm{HS,R}}\,K_{\mathrm{eff}}^{2n_{\mathrm{HS}}}
    - \Delta\Theta_{\mathrm{HS}}
  \label{eq:hotspot}
\end{align}

\noindent where $K_{\mathrm{eff}}$ is the effective load factor incorporating GIC:
\begin{equation}
  K_{\mathrm{eff}} = \sqrt{K^2 + \left(\frac{I_{\mathrm{GIC}}}{I_{\mathrm{rated}}}
    \times k_{\mathrm{gic}}\right)^2}
  \label{eq:keff}
\end{equation}
$k_{\mathrm{gic}}$ is the GIC-to-thermal-equivalent factor
(typically $0.5$--$1.0$ depending on transformer design).

\subsection{Alarm Levels}

\begin{table}[ht]
\centering
\caption{GIC protection alarm levels based on hot-spot temperature.}
\label{tab:alarms}
\begin{tabular}{lcc}
\toprule
Level & Hot-spot $\Theta_{\mathrm{HS}}$ & Action \\
\midrule
NORMAL   & $< 98^\circ$C & None \\
ELEVATED & $98$--$120^\circ$C & Operator alert, reduce load \\
CRITICAL & $120$--$140^\circ$C & Automatic load shedding \\
TRIP     & $\geq 140^\circ$C & Transformer isolation \\
\bottomrule
\end{tabular}
\end{table}

\section{Simulation Results}
\label{sec:results}

\subsection{5-Bus 400\,kV Network: GIC vs E-Field}

\begin{table}[ht]
\centering
\caption{GIC levels and alarm status vs.\ E-field intensity
  (5-bus 400\,kV, $K=1$, $\Theta_{\mathrm{amb}}=25^\circ$C).}
\label{tab:gic}
\begin{tabular}{ccccc}
\toprule
$E$ (V/km) & Max $I_{\mathrm{GIC}}$ (A) & $\Theta_{\mathrm{HS}}$ ($^\circ$C) & Alarm & Status \\
\midrule
0.5 &  28.9 & 91.0  & NORMAL   & \cmark \\
1.0 &  57.7 & 91.3  & NORMAL   & \cmark \\
2.0 & 115.5 & 92.3  & NORMAL   & \cmark \\
5.0 & 288.7 & 107.8 & ELEVATED & \cmark \\
10.0 & 577.3 & 127.0 & CRITICAL & \cmark \\
\bottomrule
\end{tabular}
\end{table}

\subsection{11-Test Pass Summary}

\begin{table}[ht]
\centering
\caption{TR-81 unit test results.}
\label{tab:tests}
\begin{tabular}{llc}
\toprule
Test & Description & Status \\
\midrule
T01 & LP model: 3-bus KCL check ($E_x$ component) & \cmark \\
T02 & LP model: 3-bus KCL check ($E_y$ component) & \cmark \\
T03 & LP model: linear E-field scaling & \cmark \\
T04 & LP model: 5-bus 400\,kV validation & \cmark \\
T05 & Thermal model: steady-state (no GIC) & \cmark \\
T06 & Thermal model: GIC-augmented $K_{\mathrm{eff}}$ & \cmark \\
T07 & Alarm: NORMAL level ($I_{\mathrm{GIC}} = 100$\,A) & \cmark \\
T08 & Alarm: ELEVATED level ($I_{\mathrm{GIC}} = 350$\,A) & \cmark \\
T09 & Alarm: CRITICAL level ($I_{\mathrm{GIC}} = 500$\,A) & \cmark \\
T10 & End-to-end pipeline: LP GIC $\to$ thermal model & \cmark \\
T11 & Euler integration accuracy ($< 5^\circ$C vs analytical) & \cmark \\
\midrule
\textbf{Total} & --- & \textbf{11/11 PASS} \\
\bottomrule
\end{tabular}
\end{table}

\begin{finding}[GIC thermal risk begins at $E > 5$\,V/km]
For the 5-bus 400\,kV network, the transformer hot-spot temperature
reaches the ELEVATED alarm threshold at $E = 5\,\text{V/km}$
($I_{\mathrm{GIC}} = 288.7\,\text{A}$) and CRITICAL at $10\,\text{V/km}$.
The 1989 Quebec event produced estimated $E \approx 10\,\text{V/km}$
at some locations, consistent with the observed transformer damage.
\end{finding}

\section{Conclusion}
\label{sec:conclusion}

\begin{finding}[LP GIC model: exact KCL verified]
The Lehtinen-Pirjola GIC solver satisfies KCL ($\sum I_{\mathrm{GIC}} = 0$)
exactly for all tested E-field configurations, confirming numerical correctness.
GIC scales linearly with E-field intensity as expected from the linear LP equations.
\end{finding}

\begin{finding}[IEC 60076-7 thermal model: 11 tests PASS]
The IEC~60076-7:2018 thermal model with GIC-augmented load factor correctly
tracks hot-spot temperature with Euler integration error $< 5^\circ$C
vs.\ the analytical solution.
\end{finding}

\begin{thebibliography}{99}
\bibitem{Gaunt2007}
C.~T.~Gaunt and G.~Coetzee,
``Transformer failures in regions incorrectly considered to have low GIC-risk,''
in \emph{Proc.\ IEEE PowerTech}, Lausanne, 2007.

\bibitem{Lehtinen1985}
M.~Lehtinen and R.~Pirjola,
``Currents produced in earthed conductor networks by geomagnetically-induced
electric fields,''
\emph{Ann.\ Geophys.}, vol.~3, no.~4, pp.~479--484, 1985.

\bibitem{IEC60076}
IEC~60076-7:2018, \emph{Power Transformers --- Part~7: Loading Guide for
Mineral-Oil-Immersed Power Transformers}, IEC, Geneva, 2018.

\bibitem{NERC2012}
NERC, \emph{Special Reliability Assessment: Effects of Geomagnetic
Disturbances on the Bulk Power System}, Atlanta, GA, 2012.

\bibitem{Boteler2014}
D.~H.~Boteler and R.~J.~Pirjola,
``Modelling geomagnetically induced currents,''
\emph{Space Weather}, vol.~15, no.~1, pp.~258--276, 2017.
\end{thebibliography}

\end{document}
